% Magnitude-sweep figure. Drop-in once analysis/aggregate_results.py has written
% magnitude_sweep.dat. This will be the paper's first figure with error bars.
%
% To use: copy magnitude_sweep.dat next to main.tex, then \input this file, or paste
% the figure body into Appendix "Two Planned Controls".
%
% The .dat has four whitespace-separated columns:
%   condition  log10_takeover  cooperate_mean  cooperate_sd
% pgfplots filters one condition per \addplot via `discard if not`.
%
% Reading the plot:
%   rising with magnitude  -> "huge number is a trap" heuristic
%   flat                   -> consistent with the bounded CARA target
%   falling                -> expected-value reasoning (the untrained baseline)
% The low end is the diagnostic: $10^3 and $10^4 are not alarming numbers.

\pgfplotsset{
  discard if not/.style 2 args={
    x filter/.code={
      \edef\tempa{\thisrow{#1}}\edef\tempb{#2}%
      \ifx\tempa\tempb\else\def\pgfmathresult{inf}\fi
    }
  }
}

\begin{figure}[t]
\centering
\begin{tikzpicture}
\begin{axis}[
    width=\linewidth,
    height=6cm,
    xlabel={Takeover payoff, $\log_{10}$ dollars},
    ylabel={Cooperate \%},
    ymin=0, ymax=100,
    xmin=2, xmax=101,
    xtick={3,4,5,6,8,11,16,25,50,100},
    xticklabel style={font=\tiny, rotate=45, anchor=north east},
    grid=major,
    grid style={dashed, gray!30},
    legend pos=outer north east,
    legend style={font=\small, draw=none},
    error bars/y dir=both,
    error bars/y explicit,
]
% One \addplot per condition. Add or remove rows to match what was actually run.
\foreach \cond/\lbl in {baseline/Baseline, sft/SFT, tie/Tie training, dpo/DPO, steering/Steering}{
  \edef\temp{\noexpand\addplot+[mark=*, discard if not={condition}{\cond}]
    table[x=log10_takeover, y=cooperate_mean, y error=cooperate_sd, col sep=space]
    {magnitude_sweep.dat};
    \noexpand\addlegendentry{\lbl}}
  \temp
}
\end{axis}
\end{tikzpicture}
\caption{Cooperate percentage against the magnitude of the Rebel takeover payoff, on 732
matched families of situations that are identical except for that payoff. Mean $\pm$ 1
standard deviation over five seeds. A rising curve would indicate a heuristic triggered by
large numbers, a flat curve is what the bounded target utility predicts, and a falling
curve is expected-value reasoning.}
\label{fig:magnitude-sweep}
\end{figure}
